% Moved from sections/materials.tex in round 44 (tables live in tables/).
% Round 75: 'band + far-IR, adopted model' column from outputs/paper1/farir_adopted_round75.json (scripts/farir_adopted_round75.py).
\begin{table*}[t]
\centering
\caption{\ClaudeText{Alternative materials through the baseline design ($z=20$~km, marginalized $3\sigma$ single-sounding threshold for the same $1$~Tg-normalized layer at the silica detection PSD; calibrated two-component background with the six sulfate parameters and the material's amplitude fitted; heritage triplet plus the $7.8$--$13~\mu$m band at $0.1~\mu$m sampling). ``Adopted model'' is the joint retrieval under the measurement model of Table~\ref{tab:error_budget}; ``scenario'' is the diagonal-floor scenario with the design's per-element floors (undisturbed calibration; post--Hunga Tonga calibration in parentheses). The uniaxial materials enter through the orientation-averaged polarizability (Sect.~\ref{sec:materials}), whose error beyond the Rayleigh limit is not quantified (Appendix~\ref{app:sensitivity}). $\lambda_\text{peak}$ and $\alpha_\text{peak}$ are the material's strongest resonance inside the band (the band-grid element nearest it) and its extinction per Tg at $20$~km (for silica the $8.80$-$\mu$m element, on which the design grid is centred). The window subset ($7.8$--$9.3~\mu$m) contains no band of the three alternatives and is blind to them ($13$--$25$~Tg).}}
\label{tab:materials}
\footnotesize
\begin{tabular}{lcccc}
\toprule
Material & $\lambda_\text{peak}$ & $\alpha_\text{peak}$ & design band, adopted model & design band, scenario \\
 & [$\mu$m] & [$10^{-7}$~m$^{-1}$\,Tg$^{-1}$] & $M_{\min}^{(3\sigma)}$ [Tg] & $M_{\min}^{(3\sigma)}$ [Tg] \\
\midrule
Silica (reference) & 8.80 & 6.2 & 0.119 & 0.101 (0.100) \\
Calcite & 11.4 & 1.0 & 0.41 & 0.80 (0.79) \\
Dolomite & 11.3 & 0.51 & 0.78 & 1.16 (1.16) \\
Alumina & 12.9 & 3.1 & 0.18 & 0.157 (0.158) \\
\bottomrule
\end{tabular}
\end{table*}

\begin{table*}[t]
\centering
\caption{\ClaudeText{The far-infrared extension of Table~\ref{tab:materials}: a $5$-$\mu$m-wide band of $0.1~\mu$m elements (range listed), with the range of the calculated per-element gas-removal floors (the median element of every band sits at the lower bound, the $\mathrm{SNR}=2000$ floor). ``Far-IR only'' is that band with the triplet and without the design band; the two combined columns add it to the design band in the diagonal-floor scenario and under the measurement model of Table~\ref{tab:error_budget}, every row in the per-shell convention of Sect.~\ref{sec:problem}. Brackets give the values for far-IR floors $3\times$ and $10\times$ the calculated ones. The far-infrared floors are calculated, not measured, and the far-infrared elements carry no empirical residual term.}}
\label{tab:materials_farir}
\footnotesize
\begin{tabular}{lcccccc}
\toprule
Material & far-IR band & far-IR floor & far-IR only & band + far-IR, scenario & band + far-IR, adopted model \\
 & [$\mu$m] & [$10^{-8}$~m$^{-1}$] & $M_{\min}^{(3\sigma)}$ [Tg] & $M_{\min}^{(3\sigma)}$ [Tg] & $M_{\min}^{(3\sigma)}$ [Tg] \\
\midrule
Silica (reference) & 18--23 & 0.34--8.1 & 0.047 [0.14, 0.47] & 0.036 [0.075, 0.095] & 0.038 [0.083, 0.113] \\
Calcite & 26--31 & 0.34--256 & 0.037 [0.11, 0.37] & 0.026 [0.078, 0.24] & 0.026 [0.076, 0.20] \\
Dolomite & 22.5--27.5 & 0.34--29 & 0.063 [0.19, 0.63] & 0.031 [0.093, 0.29] & 0.031 [0.091, 0.25] \\
Alumina & 18.25--23.25 & 0.34--6.3 & 0.10 [0.31, 1.0] & 0.071 [0.105, 0.115] & 0.075 [0.120, 0.143] \\
\bottomrule
\end{tabular}
\end{table*}
\Claude{Round 76 (T4): Table 9 split into the design-band and far-infrared tables at a readable size; the round-75 table is kept below as a comment.}

% ---- round-75 single table ----
% \begin{table*}[t]
% \centering
% \caption{Alternative materials through the baseline design and
% its far-infrared extension ($z=20$~km, marginalized $3\sigma$
% single-sounding threshold for the same $1$~Tg-normalized layer at the
% silica detection PSD; calibrated two-component background with the six sulfate parameters and the material's amplitude fitted). The ``design band, adopted model'' column is the joint retrieval under the measurement model of Table~\ref{tab:error_budget}; every other column is the diagonal-floor scenario. The uniaxial materials enter every column through the orientation-averaged polarizability (Sect.~\ref{sec:materials}). Design band: heritage triplet plus the $7.8$--$13~\mu$m
% band at $0.1~\mu$m sampling with the design's per-element floors
% (undisturbed calibration; post--Hunga Tonga calibration in parentheses);
% $\lambda_\text{peak}$ and $\alpha_\text{peak}$ are the material's
% strongest resonance inside the band (the band-grid element nearest it) and its extinction per Tg at $20$~km (for silica the $8.80$-$\mu$m element, on which the design grid is centred). Far-IR: a $5$-$\mu$m-wide band of $0.1~\mu$m elements
% (range listed), with the range of the calculated per-element
% gas-removal floors given (the median element of every band sits at
% the lower bound, the $\mathrm{SNR}=2000$ floor); ``far-IR only'' is that band with the
% triplet and without the design band, the last column is the design band plus the far-IR
% band; in these two columns the values in brackets are for
% far-IR floors $3\times$ and $10\times$ the calculated ones. The window
% subset ($7.8$--$9.3~\mu$m) contains no band of the three alternatives and
% is blind to them ($13$--$25$~Tg). The last column combines the design band under the measurement model of Table~\ref{tab:error_budget} with the far-infrared elements at the calculated floors (brackets: floors $3\times$ and $10\times$), every row in the per-shell convention of Sect.~\ref{sec:problem}; the far-infrared-only column does not involve the design band.}
% \label{tab:materials}
% \footnotesize
% \begin{adjustbox}{max width=\textwidth}
% \begin{tabular}{lccccccccc}
% \toprule
% Material & $\lambda_\text{peak}$ & $\alpha_\text{peak}$ & design band, adopted model & design band, scenario & far-IR band & far-IR floor & far-IR only & band + far-IR, scenario & band + far-IR, adopted model \\
%  & [$\mu$m] & [$10^{-7}$~m$^{-1}$\,Tg$^{-1}$] & $M_{\min}^{(3\sigma)}$ [Tg] & $M_{\min}^{(3\sigma)}$ [Tg] & [$\mu$m] & [$10^{-8}$~m$^{-1}$] & $M_{\min}^{(3\sigma)}$ [Tg] & $M_{\min}^{(3\sigma)}$ [Tg] & $M_{\min}^{(3\sigma)}$ [Tg] \\
% \midrule
% Silica (reference) & 8.80 & 6.2 & 0.119 & 0.101 (0.100) & 18--23 & 0.34--8.1 & 0.047 [0.14, 0.47] & 0.036 [0.075, 0.095] & 0.038 [0.083, 0.113] \\
% Calcite & 11.4 & 1.0 & 0.41 & 0.80 (0.79) & 26--31 & 0.34--256 & 0.037 [0.11, 0.37] & 0.026 [0.078, 0.24] & 0.026 [0.076, 0.20] \\
% Dolomite & 11.3 & 0.51 & 0.78 & 1.16 (1.16) & 22.5--27.5 & 0.34--29 & 0.063 [0.19, 0.63] & 0.031 [0.093, 0.29] & 0.031 [0.091, 0.25] \\
% Alumina & 12.9 & 3.1 & 0.18 & 0.157 (0.158) & 18.25--23.25 & 0.34--6.3 & 0.10 [0.31, 1.0] & 0.071 [0.105, 0.115] & 0.075 [0.120, 0.143] \\
% % ---- round 71: far-IR / band+far-IR columns under the ray-arithmetic mean (a), for the record: calcite 0.033 [0.099, 0.33] / 0.024 [0.072, 0.22]; dolomite 0.041 [0.12, 0.41] / 0.027 [0.080, 0.25]; alumina 0.13 [0.39, 1.3] / 0.077 [0.110, 0.119] (outputs/materials_calibrated/results.json; reproduced by materials_farir_uniaxial_round71.py to 1e-14, alumina 2e-3) ----
% % ---- round-52 rows, for the record ----
% % Silica (reference) & 8.75 & 5.6 & 0.081 (0.081) & 18--23 & 0.39--9.1 & 0.054 [0.16, 0.54] & 0.038 [0.068, 0.078] \\
% % Calcite & 11.45 & 0.81 & 0.64 (0.64) & 26--31 & 0.39--290 & 0.037 [0.11, 0.37] & 0.027 [0.081, 0.24] \\
% % Dolomite & 11.35 & 0.91 & 0.60 (0.60) & 22.5--27.5 & 0.39--33 & 0.046 [0.14, 0.46] & 0.031 [0.090, 0.26] \\
% % Alumina & 12.95 & 8.0 & 0.125 (0.124) & 18.25--23.25 & 0.39--7.1 & 0.15 [0.44, 1.5] & 0.073 [0.093, 0.098] \\
% \bottomrule
% \end{tabular}
% \end{adjustbox}
% \end{table*}
